\section{The Proposed Algorithm: MPECSS}
\label{sec:algorithm}
In this section, we present the MPECSS algorithm.
% which combines Scholtes's regularization \cite{Scholtes2001} with an adaptive path-following strategy guided by a multiplier sign test and complementarity residual to target B-stationary points of MPECs. 
We first describe the main parts of the algorithm and then present a detailed step-by-step procedure.

\subsection{Scholtes's regularization}
\label{sec:scholtes}
As discussed in Section~\ref{sec:intro}, the main source of difficulty in MPECs~\eqref{eq:general_mpec} is the complementarity conditions, which cause standard constraint qualifications to fail. In 2001, Scholtes~\cite{Scholtes2001} proposed a regularization technique, which relaxed the complementarity constraints $G_\ell(x) \cdot H_\ell(x) = 0$ to $G_\ell(x) \cdot H_\ell(x) \leq t_k$ for a sequence of positive parameters $t_k \to 0$ indexed by outer iteration $k$. This regularization transforms the original MPEC into a sequence of smooth NLP subproblems that can be solved using standard NLP solvers~\cite{FletcherLeyffer2004}. The key idea is that as $t_k$ approaches zero, the solutions of the regularized problems approach a stationary point of the original MPEC.

After replacing the complementarity constraints with their regularized approximations, the regularized problem $\text{NLP}(t_k)$ with parameter $t_k$ at outer iteration $k$ can be formulated as~\cite{Scholtes2001}:
\begin{equation}
\begin{aligned}
  \text{NLP}(t_k): \quad & \min_{x \in \mathbb{R}^n} \quad && f(x) \\
  & \text{s.t.} \quad && g_i(x) \leq 0, \quad i = 1, \ldots, m, \\
  & && h_j(x) = 0, \quad j = 1, \ldots, p, \\
  & && G_\ell(x) \ge 0, \quad \ell = 1, \ldots, q, \\
  & && H_\ell(x) \ge 0, \quad \ell = 1, \ldots, q, \\
  & && G_\ell(x) \cdot H_\ell(x) \leq t_k, \quad \ell = 1, \ldots, q.
\end{aligned}
\label{eq:nlp}
\end{equation}

The sequence of solutions $\{\text{NLP}(t_k)\}_{k=0}^\infty$ is expected to converge to a C-stationary point of the original MPEC as $t_k \to 0$~\cite{Scholtes2001}.

The regularization sequence $\{t_k\}$ is generated by a geometric reduction~\cite{Scholtes2001,HoheiselKanzowSchwartz2013}:
\begin{equation}
  t_{k+1} = \kappa \cdot t_k, \quad k = 0, 1, 2, \ldots,
  \label{eq:param_update}
\end{equation}
where $t_0 > 0$ is the initial regularization parameter and $\kappa \in (0, 1)$ is the reduction factor. As $k \to \infty$, $t_k \to 0$, bringing the regularized problems progressively closer to the original MPEC~\eqref{eq:general_mpec}. 

Each subproblem $\text{NLP}(t_{k+1})$ is initialized from the solution $x^k$ of the previous subproblem $\text{NLP}(t_k)$. This \emph{warm-start} strategy means the solver follows a continuous \emph{regularization path} $\{x^k\}$ as $t_k$ decreases, rather than solving each subproblem from scratch~\cite{RalphWright2004,HoheiselKanzowSchwartz2013}.

A solution $x^k$ of $\text{NLP}(t_k)$ is accepted when the KKT conditions are satisfied to a tolerance $\varepsilon_k > 0$. In practice, $\varepsilon_k$ is reduced alongside $t_k$ (e.g., $\varepsilon_k = O(t_k)$; see~\cite{HoheiselKanzowSchwartz2013,KanzowSchwartz2015}) to ensure that the accuracy of the solutions improves as $t_k$ gets smaller.

Scholtes~\cite{Scholtes2001} proved that limit points of the regularized sequence are at best C-stationary for the MPEC. In general, C-stationarity is \emph{strictly weaker} than B-stationary
 and the standard method provides no mechanism to detect or target B-stationary points~\cite{Scholtes2001,HoheiselKanzowSchwartz2013}. Alternative relaxation formulations have been proposed that guarantee stronger stationarity under weaker
 conditions (e.g., Kanzow and Schwartz~\cite{KanzowSchwartz2013} achieve
 M-stationarity under MPEC-LICQ; DeMiguel et al.~\cite{DeMiguel2005} propose a two-sided relaxation), but their theoretical stationarity guarantees can
 fail when NLP solvers return approximate solutions~\cite{KanzowSchwartz2015}. MPECSS
 takes a different approach: it preserves Scholtes's product relaxation which is
 numerically robust and fills the stationarity gap by combining the regularization path
 with an \emph{adaptive path-following strategy} guided by
the complementarity residual's improvement rate. At each
outer iteration, the \emph{multiplier sign test} (Section~\ref{sec:stationarity})
checks the $G$-multipliers and $H$-multipliers of
$\text{NLP}(t_k)$ to confirm (or exclude) S-stationarity of the
current iterate, successfully screening high-quality candidates. The full description of this adaptive mechanism is given
in Section~\ref{sec:adaptive_strategy}. Notably, S-stationarity implies B-stationarity by
Scheel and Scholtes~\citep{ScheelScholtes2000}, and under the
MPEC-LICQ constraint qualification, S-stationarity and B-stationarity
are equivalent. In the implementation, however, a sign-test pass is
carried into Phase~III, where B-stationarity is certified either by
LICQ equivalence or by the LPEC check when LICQ is absent or
fails to hold.

\textit{Note:} MPECSS's code accepts benchmark models whose native complementarity format is
not exactly the textbook form~\eqref{eq:nlp}. For MPECLib and
NOSBENCH, finite lower bounds on the complementarity functions are
shifted into internal variables
\[
  \widehat G_\ell(x) := G_\ell(x)-\ell^G_\ell,\qquad
  \widehat H_\ell(x) := H_\ell(x)-\ell^H_\ell,
\]
so that the solver always works with nonnegative internal pairs
$\widehat G_\ell,\widehat H_\ell \ge 0$. For MacMPEC, the input handler also
supports mixed and box-complementarity structures, including free
$G$-components and upper-bound pairs of the form
$(-G_\ell(x))(u_\ell-H_\ell(x)) \le t_k$. The notation in the sequel is
therefore to be understood as the solver's internal complementarity
representation after this benchmark-specific normalization.


\subsection{Phase~I: Feasibility}
\label{sec:feasibility}

The main MPECSS loop (Section~\ref{sec:adaptive_strategy}) requires a starting point that approximately satisfies the complementarity conditions $G_\ell(x)^{\top}\, H_\ell(x) \approx 0$ with $G_\ell(x) \geq 0$ and $H_\ell(x) \geq 0$. A user-provided initial point~$x_0$ may be far from complementarity-feasible (i.e., the products $G_\ell(x_0)^{\top}\, H_\ell(x_0)$ may be large). To address this, MPECSS begins with an optional \emph{feasibility phase} (Phase~I), following standard MPEC feasibility preprocessing approaches~\cite{BaumruckerBiegler2008,NurkanovicLeyffer2025}, that finds a good starting point satisfying the complementarity conditions. Before solving, the initial point~$x_0$ is pushed into the strict interior of the variable bounds by a fraction $\nu = 0.1$, i.e., for each dimension~$i$ with finite bounds $[\ell_i, u_i]$, the starting value is restricted to $[\ell_i + \nu(u_i - \ell_i),\; u_i - \nu(u_i - \ell_i)]$. 
% This \emph{interior push} avoids Jacobian singularities that interior-point methods can encounter when initialized exactly on variable bounds~\cite{WachterBiegler2006}.

Phase~I then solves a sequence of up to three progressively more robust NLP formulations to find a complementarity-feasible point. The general structure minimizes $\Theta(x)$ of the complementarity conditions, subject to the original MPEC feasibility constraints:
\begin{equation}
\begin{aligned}
  & \min_{x \in \mathbb{R}^n} \quad && \Theta(x) \\
  & \text{s.t.} \quad && g_i(x) \leq 0, \quad i = 1, \ldots, m, \\
  & && h_j(x) = 0, \quad j = 1, \ldots, p, \\
  & && G_\ell(x) \geq 0, \quad \ell = 1, \ldots, q, \\
  & && H_\ell(x) \geq 0, \quad \ell = 1, \ldots, q.
\end{aligned}
\label{eq:phase_i}
\end{equation}
The formulation of the objective $\Theta(x)$ is progressively adjusted in a cascade if the previous attempt fails to achieve a satisfactory complementarity residual ($r < 10^{-6}$). These attempts correspond to: (1) \emph{L2 product minimisation} $\Theta(x) = \sum_{\ell=1}^q \big(G_\ell(x) H_\ell(x)\big)^2$ to standardly push products to zero~\cite{NocedalWright2006}; (2) \emph{Smooth L1 minimisation} $\Theta(x) = \sum_{\ell=1}^q \sqrt{\big(G_\ell(x) H_\ell(x)\big)^2 + \epsilon}$ (with $\epsilon = 10^{-10}$) to provide a less aggressive gradient when ill-conditioned~\cite{Fischer1992,FacchineiPang2003}; and (3) an \emph{Epigraph formulation} which solves $\min_{x, \tau} \tau$ subject to the constraints in~\eqref{eq:phase_i} as well as $\big(G_\ell(x) H_\ell(x)\big)^2 \leq \tau$ for all $\ell = 1, \ldots, q$ and $\tau \geq 0$, achieving stable gradients on severel infeasible problems~\cite{RockafellarWets1998}. For box-constrained complementarity pairs, the formulation aggregates both the lower residual $G_\ell(x) H_\ell(x)$ and, when a finite upper bound $u_\ell$ is present, the upper-bound residual $G_\ell(x)(u_\ell - H_\ell(x))$.

No regularization parameter appears in any of these formulations; each is a standard smooth NLP solved by IPOPT~\cite{WachterBiegler2006}. Upon convergence of any formulation, the resulting point $\bar{x}$ satisfies $G_\ell(\bar{x}) \cdot H_\ell(\bar{x}) \approx 0$ for all~$\ell$ and lies in the feasible region of the original constraints. This point is then used as the warm-start~$x^0$ for the main Scholtes outer loop. Algorithm~\ref{alg:phase1} summarizes the practical Phase~I procedure used to generate a reliable warm start for the adaptive outer loop.

\begin{algorithm}[H]
\caption{Phase~I Feasibility}
\label{alg:phase1}
\footnotesize
\begin{algorithmic}[1]
\Require Problem data $P$ with complementarity functions and bounds, initial point $x_0$, solver options, maximum attempts $a_{\max}$, restart budget $r_{\max}$, seed $s$
\Ensure Warm-start point $x^0$ and Phase~I metrics
\State $n_x \gets \dim(x_0)$
\State $r_0 \gets \mathrm{comp\_res}(x_0, P)$
\If{$n_x > 3000$}
  \State \Return $x^0 \gets x_0$
\EndIf
\State $(\ell_x, u_x) \gets$ variable bounds from $P$
\State $x_0 \gets \textsc{InteriorPush}(x_0, \ell_x, u_x, 0.1)$
\State $x_{\mathrm{best}} \gets x_0$, $r_{\mathrm{best}} \gets r_0$
\State $x_{\mathrm{warm}} \gets x_0$, $r_{\mathrm{warm}} \gets r_0$
\State $\sigma \gets \max(1, \|x_0\|_2)$, $D \gets 50$
\State Initialise residual and objective metrics
\For{$a = 0, \ldots, a_{\max} - 1$}
  \State Solve Phase~I NLP regime $a$ from $x_0$ if $a=0$, else from $x_{\mathrm{warm}}$
  \State Obtain candidate $\bar{x}$ and residual $r$
  \State $d \gets \|\bar{x} - x_0\|_2 / \sigma$
  \If{$r < r_{\mathrm{warm}}$}
    \State $x_{\mathrm{warm}} \gets \bar{x}$, $r_{\mathrm{warm}} \gets r$
  \EndIf
  \If{$r < r_{\mathrm{best}}$ \textbf{and} $(d < D \textbf{ or } r < 10^{-6})$}
    \State $x_{\mathrm{best}} \gets \bar{x}$, $r_{\mathrm{best}} \gets r$
  \EndIf
  \If{$r < 10^{-6}$}
    \State \textbf{break}
  \EndIf
\EndFor
\If{$r_{\mathrm{best}} > 10^{-4}$ \textbf{and} $r_{\max} > 0$}
  \State Build deterministic restart candidates from lower bounds, upper bounds, midpoint, and Gaussian perturbations using a seed mixed with the problem name
  \For{$j = 0, \ldots, \min(r_{\max}, |\mathcal C|) - 1$}
    \State $x_{\mathrm{cand}} \gets \mathcal C_j$
    \For{$a = 0, 1$}
      \State Solve Phase~I NLP regime $a$ from $x_{\mathrm{cand}}$ if $a=0$, else from $x_{\mathrm{best}}$
      \State Obtain candidate $\bar{x}$ and residual $r$
      \State $d \gets \|\bar{x} - x_0\|_2 / \sigma$
      \If{$r < r_{\mathrm{best}}$ \textbf{and} $(d < D \textbf{ or } r < 10^{-6})$}
        \State $x_{\mathrm{best}} \gets \bar{x}$, $r_{\mathrm{best}} \gets r$
      \EndIf
      \If{$r_{\mathrm{best}} < 10^{-4}$}
        \State \textbf{break}
      \EndIf
    \EndFor
    \If{$r_{\mathrm{best}} < 10^{-4}$}
      \State \textbf{break}
    \EndIf
  \EndFor
\EndIf
\State \Return $x^0 \gets x_{\mathrm{best}}$ and associated metrics
\end{algorithmic}
\end{algorithm}

Three baseline safety mechanisms enhance Phase~I's reliability. First,
a \textit{structured multistart} samples deterministic lower-bound,
upper-bound, and midpoint points together with Gaussian-perturbed
starts whenever the best complementarity residual remains above
$10^{-4}$. Second, a \textit{displacement guard} accepts a candidate
only if its normalized displacement from the initial point is below a
fixed limit of $50$, unless the residual is already below $10^{-6}$.
Third, the implementation records Phase~I objective values for
diagnostics, but the acceptance rule is driven by complementarity-residual
improvement together with the displacement guard rather than by a
separate original-objective filter.

In practice, the solver may fail to find an initial search direction from the user-provided starting point and report infeasibility immediately. To improve startup robustness without changing the theoretical Scholtes framework~\cite{Scholtes2001}, MPECSS cycles through the three objective formulations described above. Each formulation offers a different objective geometry, and the cascade is arranged so that a more robust formulation is tried whenever the preceding one encounters numerical difficulty. For very large problems ($n > 3000$), Phase~I is skipped entirely as it becomes computationally expensive; in this case, the user-provided starting point is passed directly to Phase~II.

When Phase~I achieves a complementarity residual $r_0 \leq 100\,
\varepsilon_{\mathrm{tol}}$, the implementation performs an early
stationarity evaluation before entering the Phase~II outer loop: it
solves one NLP at a small regularization parameter
$\bar{t} = \max(0.1\,r_0,\, 10^{-12})$ to extract multipliers and runs
the sign test.  Three outcomes are identified:
If the sign test passes and $r_0 \leq \varepsilon_{\mathrm{tol}}$,
Phase~II is skipped and the point proceeds directly to Phase~III.
If $r_0 \leq \varepsilon_{\mathrm{tol}}$ but the sign test fails
and the biactive count satisfies
$|\mathcal{I}_{00}| \leq N_{\mathrm{LPEC}}$ (default $N_{\mathrm{LPEC}}=15$),
Phase~III is executed directly since the LPEC enumeration
($2^{|\mathcal{I}_{00}|}$ branches) is computationally manageable.
Otherwise, Phase~II runs normally.
Additionally, when Phase~I produces $r_0 < t_0$, the initial
regularization parameter is fast-forwarded to
$t_0 \gets \max(10\,r_0,\, \varepsilon_{\mathrm{tol}}\,\tau)$ to
preserve the refined starting point and avoid redundant outer
iterations at large $t$ values.

\subsection{Phase II: Adaptive Path-Following Strategy}
\label{sec:adaptive_strategy}
The main idea in MPECSS's Phase~II is to integrate a robust residual-based 
path-following framework to guide the sequence $\{t_k\}$, while applying a
\emph{multiplier sign test} to screen and isolate S-stationary candidates~\cite{ScheelScholtes2000}. 

\subsubsection{Solve regularized NLP subproblems}
\label{sec:solveRegNLPs}
Phase~I provides a robust warm start for the main loop. At each outer
iteration $k$, MPECSS builds the smooth subproblem $\text{NLP}(t_k)$
in CasADi and solves it with IPOPT~\cite{WachterBiegler2006} as its
primary NLP solver. In the reported benchmark runs, IPOPT uses the
open-source MUMPS linear solver~\cite{Amestoy2001} with tuned settings
for robustness on ill-conditioned MPEC subproblems.\footnote{%
  MUMPS is configured with increased workspace allocation
  ($\texttt{mumps\_mem\_percent}=500$), Hungarian scaling
  ($\texttt{mumps\_scaling}=77$), and pivot tolerance
  $\texttt{mumps\_pivtol}=10^{-6}$. IPOPT uses gradient-based NLP
  scaling, up to 4 second-order correction steps
  ($\texttt{max\_soc}=4$), and 3 watchdog trial iterations.}
The codebase also provides an optional small-problem
SQP$+$qpOASES~\cite{Ferreau2014qpOASES} for problems with
$n \leq 400$ variables, when qpOASES is available. %In the current
% implementation and archived benchmark runs, interior-point methods are
% used as the default because they were empirically the most reliable
% option across the tested problem families.
For larger problems, CasADi's MX symbolic type is used together with a reverse-mode-favoring automatic differentiation setup. Solver templates and Jacobian functions are managed through bounded LRU caches with automatic removal, which helps keep memory usage predictable during long benchmark runs. Algorithm~\ref{alg:phase2-main} provides the structure framework of Phase~II.

\begin{algorithm}[htbp]
\caption{Phase~II Main Adaptive Path-Following Loop}
\label{alg:phase2-main}
\begin{algorithmic}[1]
\Require Warm-start $x^0$, initial regularization $t_0$, reduction factor $\kappa$, tolerance $\varepsilon_{\mathrm{tol}}$, maximum iterations $k_{\max}$
\Ensure Candidate point $x^*$ for Phase~III
\State Set $k \gets 0$; initialise history buffers, stagnation counters, and best-point tracker $\mathrm{best}$
\While{$k < k_{\max}$}
  \State Solve the regularized NLP at $t_k$ from warm-start $x^k$ to obtain $(x^{k+1},\lambda^{k+1})$
  \State Convert multipliers and compute the homotopy residual $r_{k+1}$
  \State Run the sign test on $\mathcal{I}_{00}(x^{k+1})$ and update $\mathrm{best}$ if the candidate improves the incumbent
  \If{$r_{k+1} \le \varepsilon_{\mathrm{tol}}$ \textbf{and} sign test passes}
    \State Set $x^* \gets x^{k+1}$ and \textbf{break}
  \EndIf

  \State Update $t_{k+1}$ via Algorithm~\ref{alg:phase2-update}
  \State $k \gets k+1$
\EndWhile
\If{no convergence}
  \State $x^* \gets \mathrm{best}$
\EndIf
\State \Return $x^*$
\end{algorithmic}
\end{algorithm}

When the primary solve fails (e.g., due to infeasibility detection
or numerical errors), MPECSS uses a small sequence of IPOPT retry configurations, including a Mehrotra-style predictor-corrector variant, relaxed tolerances, and gradient-based scaling; if these attempts fail, it can switch to a Fischer--Burmeister reformulation of the complementarity constraints~\cite{WachterBiegler2006,Fischer1992,Kanzow1996}.

Each subproblem $\text{NLP}(t_{k+1})$ is warm-started from the solution
$x^k$ of the previous subproblem $\text{NLP}(t_k)$, including primal
and dual variables.\footnote{%
  IPOPT warm-start parameters: $\texttt{warm\_start\_bound\_push} =
  \texttt{warm\_start\_bound\_frac} = \texttt{warm\_start\_slack\_bound\_frac}
  = \texttt{warm\_start\_slack\_bound\_push} = 10^{-9}$.}
This \emph{warm-start} strategy means the solver follows a continuous
\emph{regularization path} $\{x^k\}$ as $t_k$ decreases, rather than
solving each subproblem from scratch, greatly reducing the total number
of NLP iterations.


\subsubsection{Constraint formulation and multiplier sign test}
\label{sec:stacking_signs}
For implementation consistency, the regularized constraints are passed to the NLP solver in a fixed order:
\begin{equation}
g(x)=\begin{bmatrix}
g_{\mathrm{orig}}(x)\\
G(x)+\delta_k\\
H(x)+\delta_k\\
G(x)\circ H(x)-t_k
\end{bmatrix},
\label{eq:stacked_constraints}
\end{equation}
where the tag $\mathrm{orig}$ means ``original'' and
$g_{\mathrm{orig}}(x)$ collects the inequality constraints
inherited from the original MPEC model (before adding regularization
terms), $\circ$ denotes componentwise product, and $\delta_k \geq 0$
is an optional regularization shift.\footnote{%
  In all reported numerical experiments, $\delta_k = 0$; the shift is
  retained in the formulation for generality.} Let
$\lambda^{\mathrm{cas}}$ be the solver multiplier vector associated
with~\eqref{eq:stacked_constraints}, where the tag ``$\mathrm{cas}$''
means multipliers returned in the CasADi/IPOPT solver format.
Because CasADi returns non-positive multipliers at active lower bounds
for the $G$-constraints and $H$-constraints, MPECSS converts these complementarity
constraints to MPEC-oriented multipliers by negation:
\begin{equation}
\gamma^{G,k} = -\lambda^{\mathrm{cas}}_{G},\qquad
\gamma^{H,k} = -\lambda^{\mathrm{cas}}_{H},
\label{eq:multiplier_conversion}
\end{equation}
The multiplier of the relaxed product constraint is still extracted and
stored for test purposes, but it is \emph{not} part of the
S-stationarity sign test. All sign tests in MPECSS are therefore
performed only on the converted multipliers in
Equation~\eqref{eq:multiplier_conversion}.

At each outer iteration $k$, after solving $\text{NLP}(t_k)$, we
extract the converted multipliers $(\gamma^{G,k},\gamma^{H,k})$ for
the shifted $G$-constraints and $H$-constraints. For biactive indices
$\ell \in \mathcal{I}_{00}(x^k)$, the implemented sign test checks
\[
\gamma^{G,k}_\ell \ge -\tau,\qquad
\gamma^{H,k}_\ell \ge -\tau,
\]
with a tolerance chosen in the implementation as
$\tau=10^{-6}$ by default. The adaptive quantity
$\sigma_k=\max\{10^{-8},\sqrt{t_k}\}$ is used instead for detecting the
numerically biactive set. If the sign test passes and the complementarity
residual is sufficiently small, Phase~II returns an S-stationarity
candidate to Phase~III, where final B-stationarity is certified by the
implemented LICQ/LPEC logic.

\subsubsection{Adaptive $t$-update strategy}
\label{sec:adaptive_t}
In the standard Scholtes regularization, the regularization parameter is
reduced by a fixed geometric rule $t_{k+1} = \kappa \, t_k$~\cite{Scholtes2001}
regardless of the solver's progress. MPECSS replaces this with an
\emph{adaptive update strategy} that monitors the
complementarity residual
\begin{equation}
  r_k \;=\; \max_{\ell = 1,\ldots,q} \bigl|G_\ell(x^k)\,H_\ell(x^k)\bigr|
  \label{eq:comp_res}
\end{equation}
at every outer iteration $k$ and adjusts $t_k$ accordingly. For
stationarity screening only, the code also tracks the biactive
residual
\begin{equation}
  \eta_k \;=\; \max_{\ell = 1,\ldots,q} \;\min\!\bigl(|G_\ell(x^k)|,\; |H_\ell(x^k)|\bigr).
  \label{eq:biactive_res}
\end{equation}
Similar to the improvement ratio used in trust-region methods~\cite{NocedalWright2006},
we define the relative improvement
\begin{equation}
  \rho_k \;=\; \frac{r_{k-1} - r_k}{\max(r_{k-1},\, \epsilon_{\mathrm{floor}})},
  \label{eq:rho}
\end{equation}
where $\epsilon_{\mathrm{floor}}$ is a small positive guard to prevent division by zero (using the smallest normal float in the implementation). A value $\rho_k \approx 1$ indicates that the complementarity residual was nearly eliminated in one step, while small positive $\rho_k$ indicates stagnation. The algorithm monitors whether convergence is sustained and the biactive set has stabilised. Specifically, define the \emph{pathing condition} as
\begin{equation}
  \alpha_L \;\le\; \rho_k \;\le\; \alpha_U,
  \label{eq:pathing}
\end{equation}
with $\alpha_L = 0.2$ and $\alpha_U = 5.0$, where $\rho_k$ is the relative improvement~\eqref{eq:rho}. The solver is said to be in \emph{pathing mode} at iteration $k$ when~\eqref{eq:pathing} holds, the biactive count $|\mathcal{I}_{00}|$ has been constant over the last three iterates, and $k \ge 4$. Pathing mode is \emph{sustained} when the pathing counter $c \ge 2$ (i.e., the condition has held for at least two consecutive iterations). 

\text{Path 1:} When pathing mode is sustained and $\rho_k > 0.5$---indicating strong, steady convergence---the algorithm uses an extra-aggressive reduction:
\begin{equation}
  t_{k+1} = \kappa^2 \, t_k.
  \label{eq:superlinear}
\end{equation}

\text{Path 2:}
When pathing mode is sustained and $\rho_k > 0.1$, the algorithm maintains momentum with a capped reduction:
\begin{equation}
  t_{k+1} = \min(\kappa,\,0.5)\, t_k.
  \label{eq:fast}
\end{equation}

The cap at $0.5$ prevents overly aggressive reduction when $\kappa$ is set close to~$1$, ensuring robustness across a wide range of user-supplied $\kappa$ values.


\text{Path 3:}
Stagnation is tracked by a counter $s$ that increments whenever $\rho_k < 0.05$ and resets to zero otherwise. Two sub-paths handle escalating stagnation. When $s \ge 4$, an aggressive jump is executed and the counter is reset:

\begin{equation}
   t_{k+1} = \kappa^2\, t_k, \qquad s \gets 0 \qquad \text{(Path 3a)}.
  \label{eq:stagnation_jump}
\end{equation}

With the default $\kappa = 0.5$ this performs an aggressive reduction ($t_{k+1} = 0.25 t_k$), forcing the solver to move past the stagnation point. When $2 \le s < 4$, a milder post-stagnation reduction is used:

\begin{equation}
  t_{k+1} = \min(\kappa,\,0.5)\, t_k \qquad \text{(Path 3b)}.
  \label{eq:post_stagnation}
\end{equation}

\text{Path 4:}
In all other cases---pathing mode is not sustained and the stagnation count is below threshold---the algorithm falls back to the standard reduction:
\begin{equation}
  t_{k+1} = \kappa\, t_k.
  \label{eq:slow}
\end{equation}

% An alternative to the implemented update is to set $t_{k+1} = r_k$
% directly,
% inspired by the theoretical sensitivities noted
% in~\cite{HoheiselKanzowSchwartz2013}. This update can converge faster
% but is slightly less robust on problems with oscillatory complementarity
% residuals. In practice, we use the adaptive rule above as the default.

Algorithm~\ref{alg:phase2-update} summarizes the implemented adaptive
update logic.

\begin{algorithm}[H]
\caption{Phase~II Adaptive $t$-Update}
\label{alg:phase2-update}
\begin{algorithmic}[1]
\Require Current $t_k$, residuals $(r_k,r_{k-1})$, stagnation counter $s$, pathing counter $c$, biactive count history, factors $(\kappa,\alpha_L,\alpha_U)$, iteration index $k$, remaining jump budget $b$
\Ensure Next regularization parameter $t_{k+1}$, updated counters $(s,c,b)$
\State Compute relative improvement $\rho_k$ from~\eqref{eq:rho}
\State Update stagnation counter: \textbf{if} $\rho_k < 0.05$ \textbf{then} $s \gets s+1$ \textbf{else} $s \gets 0$
\State Set \textit{biactive\_stable} $\gets$ (last 3 iterates share the same $|\mathcal{I}_{00}|$)
\State Evaluate pathing condition~\eqref{eq:pathing}: \textit{cond} $\gets$ ($k \ge 4$ \textbf{and} \textit{biactive\_stable} \textbf{and} $\alpha_L \le \rho_k \le \alpha_U$)
\State Update pathing counter: \textbf{if} \textit{cond} \textbf{then} $c \gets c+1$ \textbf{else} $c \gets \max(0,c-1)$
\State Set \textit{pathing} $\gets$ ($c \ge 2$)
\If{\textit{pathing} \textbf{and} $\rho_k > 0.5$}
  \State \Comment{Path 1: superlinear mode}
  \State $t_{k+1} \gets \kappa^2\, t_k$
\ElsIf{\textit{pathing} \textbf{and} $\rho_k > 0.1$}
  \State \Comment{Path 2: fast geometric mode}
  \State $t_{k+1} \gets \min(\kappa,\,0.5)\,t_k$
\ElsIf{$s \ge 4$}
  \State \Comment{Path 3a: aggressive stagnation jump; reset $s \gets 0$}
  \State \textbf{if} $b>0$ \textbf{then} $t_{k+1} \gets \kappa^2\, t_k$, $b \gets b-1$, $s \gets 0$ \textbf{else} $t_{k+1} \gets \kappa\, t_k$
\ElsIf{$s \ge 2$}
  \State \Comment{Path 3b: post-stagnation fast}
  \State $t_{k+1} \gets \min(\kappa,\,0.5)\,t_k$
\Else
  \State \Comment{Path 4: conservative mode}
  \State $t_{k+1} \gets \kappa\, t_k$
\EndIf
  \State $t_{k+1} \gets \max(10^{-14},\, t_{k+1})$
\State \Return $t_{k+1}$
\end{algorithmic}
\end{algorithm}

In addition to the core $t$-update, the implementation includes several
practical safeguards for difficult problems. After Phase~I, if the
warm-start is already nearly complementarity-feasible, the code can
bypass Phase~II and proceed directly to Phase~III; alternatively, it can
fast-forward the initial $t_k$. During Phase~II, every outer NLP
solve is protected by a wall-clock guard, and solver failures trigger
the retry configurations described in
Section~\ref{sec:solveRegNLPs} before the current iterate is discarded.

The aggressive stagnation jump (Path~3a in Algorithm~\ref{alg:phase2-update})
is governed by a hard cap of $b = 500$ total jump triggers; exceeding
this cap prevents unbounded adaptive cycling and the outer loop
declares stagnation. Additionally, the implementation monitors an
\emph{outer stagnation guard}; if the complementarity residual changes
by less than $0.1\%$ over a sliding window of 25~iterations, the outer
loop exits early. 

\subsection{Phase III: B-Stationarity Certification}
\label{sec:phase_iii}
After Phase~II terminates with a candidate $x^*$, MPECSS enters a
certification and refinement stage inspired by MPECopt-style active-set
refinement and directional certification~\cite{NurkanovicLeyffer2025}.
In the core solver, this stage first performs the LICQ/LPEC logic on
the best complementarity-feasible point. In the benchmark pipeline used
for the numerical experiments, it is then followed by BNLP polishing, an
LPEC refinement loop, and a final post-check for B-stationarity.

When Phase~II exits with $\varepsilon_{\mathrm{tol}} < r_{\mathrm{best}}
\leq 1000\,\varepsilon_{\mathrm{tol}}$, the benchmark pipeline performs
a final push before certification. Gentle $t$ reduction solves
$\mathrm{NLP}(t)$ at progressively smaller $t$ values (factors $0.5$,
$0.1$, $0.01$ of the current estimate), rejecting any solution whose
iterate norm diverges by more than $10\times$ from the current point.
BNLP polish tries the BNLP~\eqref{eq:phase_i} at multiple active-set
tolerances and accepts only if it improves $r_{\mathrm{comp}}$.
Partition flipping flips, for each biactive index, the active-set
assignment $G_\ell = 0 \leftrightarrow H_\ell = 0$ and re-solves the
BNLP to find a feasible partition. If all strategies fail but
$r_{\mathrm{best}} \leq 10\,\varepsilon_{\mathrm{tol}}$, the point is accepted under a
relaxed tolerance and classified as C-stationary. This three-strategy
cascade recovers problems whose residual is only slightly
above the target tolerance. Algorithm~\ref{alg:phase3} summarizes the reported certification and
refinement sequence.

\begin{algorithm}[H]
\caption{Phase~III B-Stationarity Certification and Refinement}
\label{alg:phase3}
\begin{algorithmic}[1]
\Require Candidate point $x^*$ from Phase~II
\Ensure Final status and refined point
\If{$x^*$ is complementarity-feasible}
  \State Check MPEC-LICQ at $x^*$
  \If{sign test passes \textbf{and} MPEC-LICQ holds}
    \State Certify $x^*$ as B-stationary by S/B equivalence
  \Else
    \State Run LPEC-based certification on $x^*$
  \EndIf
\EndIf
\State Run BNLP polishing; accept only if complementarity remains small and the objective does not worsen materially
\State Run trust-region LPEC refinement and a final B-stationarity post-check
\State Return the final point and certification status
\end{algorithmic}
\end{algorithm}

If the final iterate passes the multiplier sign test and MPEC-LICQ holds at that point, then the implementation certifies B-stationarity by the equivalence between S- and B-stationarity under MPEC-LICQ.

Even a good solution can be numerically rough. We solve a final simplified problem (BNLP) by fixing an active-set partition. Given index sets $\mathcal{I}_1 = \{\ell : G_\ell(x^*) \approx 0\}$ and $\mathcal{I}_2 = \{\ell : H_\ell(x^*) \approx 0\}$, the BNLP subproblem is:
\begin{equation}
\begin{aligned}
  \text{BNLP}(x^*,\mathcal{I}_1,\mathcal{I}_2): \quad & \min_{x \in \mathbb{R}^n} \; f(x) \\
  & \text{s.t.} \quad g(x) \leq 0, \; h(x) = 0, \\
  & \quad\; G_\ell(x) = 0 \; (\ell \in \mathcal{I}_1), \; H_\ell(x) \geq 0 \; (\ell \in \mathcal{I}_1), \\
  & \quad\; H_\ell(x) = 0 \; (\ell \in \mathcal{I}_2), \; G_\ell(x) \geq 0 \; (\ell \in \mathcal{I}_2).
\end{aligned}
\label{eq:bnlp}
\end{equation}
For biactive indices ($\ell \in \mathcal{I}_{00}$), either $G_\ell$ or
$H_\ell$ must be assigned to the equality branch; the implementation
first uses the primary partition and then tries limited single-flip
alternatives. In the core solver, a polished point is accepted only if
its complementarity residual remains within the configured tolerance and
its objective value does not worsen by more than a small
relative/absolute tolerance. In the benchmark harness, an additional
post-check may still relabel near-feasible failures with
$r_{\mathrm{comp}} \le 10\,\varepsilon_{\mathrm{tol}}$ after the final
B-stationarity check.

If sign test plus LICQ
do not already certify the point, we use the geometric LPEC check to
determine whether a feasible descent direction exists~\cite{LuoPangRalph1996,FletcherLeyfferRalphScholtes2006}.
The implemented LP model includes the active original constraints,
active complementarity-side equalities, bound-consistent trust-region
limits $|d_i| \le \rho$, and branch-dependent inequalities for both
lower-bound and upper-bound biactive complementarity indices. For a
point with $|\mathcal{I}_{00}|$ biactive pairs, the implementation
enumerates all XOR branches (either the $G_\ell$ side or the $H_\ell$
side stays active per biactive index), capped at $2^{15}$ branches for
tractability, and solves the resulting LPs with SciPy's
\texttt{linprog(..., method='highs')} backend~\cite{HuangfuHall2018}.
If the branch enumeration times out or hits the branch cap without
finding a descent direction, the outcome is treated as
\texttt{stationarity\_unverifiable} rather than as a certificate. In
benchmark mode, this certification is followed by a trust-region
refinement loop with parameters
$\rho_{\mathrm{init}} = 10^{-2}$, $\gamma_L = 0.1$,
$N_{\mathrm{out}} = 30$, $N_{\mathrm{in}} = 15$,
$tol_B = 10^{-10}$, and $tol_{\mathrm{comp}} = 10^{-8}$.

In the reported experiments, every problem is post-processed by the full
benchmark pipeline: the Phase~II output is passed through BNLP
polishing, LPEC refinement, and a final B-stationarity post-check,
even when the initial sign test is favorable. This benchmark pipelineis
therefore slightly richer than a direct call to \texttt{run\_mpecss()},
and the numerical section reports the output of that pipeline rather than
the raw Phase~II endpoint alone.


% =========================================================================
